%\documentclass[conference]{IEEEtran}
%\documentclass[onecolumn]{IEEEtran}
%\documentclass[twocolumn]{IEEEtran}
\documentclass[journal]{IEEEtran}
% Some Computer Society conferences also require the compsoc mode option,
% but others use the standard conference format.
%
% If IEEEtran.cls has not been installed into the LaTeX system files,
% manually specify the path to it like:
% \documentclass[conference]{../sty/IEEEtran}

% Some very useful LaTeX packages include:
% (uncomment the ones you want to load)

% *** MISC UTILITY PACKAGES ***
%
%\usepackage{ifpdf}
% Heiko Oberdiek's ifpdf.sty is very useful if you need conditional
% compilation based on whether the output is pdf or dvi.
% usage:
% \ifpdf
%   % pdf code
% \else
%   % dvi code
% \fi
% The latest version of ifpdf.sty can be obtained from:
% http://www.ctan.org/pkg/ifpdf
% Also, note that IEEEtran.cls V1.7 and later provides a builtin
% \ifCLASSINFOpdf conditional that works the same way.
% When switching from latex to pdflatex and vice-versa, the compiler may
% have to be run twice to clear warning/error messages.

\usepackage{amsmath}
\usepackage{amsfonts}
\usepackage{amssymb}
\usepackage{mathrsfs}
\usepackage{booktabs}
\usepackage{bm}
\usepackage{graphicx}
\usepackage{subfigure}
\usepackage{color}
\usepackage{multirow}
\usepackage{pifont}
%\usepackage[title]{appendix}

\newtheorem{definition}{Definition}
\newtheorem{proposition}{Proposition}
\newtheorem{assumption}{Assumption}
\newtheorem{theorem}{Theorem}
\newtheorem{proof}{Proof}
\newtheorem{lemma}{Lemma}
\newtheorem{remark}{Remark}
\newtheorem{corollary}{Corollary}


\usepackage[justification=centering]{caption}

\ifCLASSINFOpdf

\else

\fi


\hyphenation{op-tical net-works semi-conduc-tor}

\begin{document}

\title{A Privacy-Preserving Distributed Secondary Control Framework for Islanded AC Microgrids: Local Physical Analysis and Public-History Indistinguishability}

\author{Jiashun Lu
\thanks
{%This work was supported by NSFC 62473204.   
%Postgraduate Research \& Practice Innovation Program of Jiangsu Province under Grant SJCX25\_0353.
%T. Zhang, Y. Liu and Y. Yang 
% The authors are with Nanjing University of Posts and Telecommunications. (Corresponding author: Yang Yang, e-mail: yyang@njupt.edu.cn.)
%, Nanjing 210023, China.  %the College of Automation, 
}
}
\maketitle

\begin{abstract}
Distributed secondary control in islanded AC microgrids must coordinate physical regulation with communication constraints and incomplete adversary access. This paper studies a frozen secondary-control architecture equipped with a public/private virtual coordination decomposition and a passive public-history observation model. The analysis is deliberately restricted to a local-before-exit theorem boundary. First, we establish local well-posedness and compact-region comparison properties for the physical closed loop, including the prescribed-performance coordinates, the finite privacy-residual bound, and actuator feasibility on a verified bootstrap design region. Second, we construct a non-nominal private realization that produces the same complete public history as a nominal realization under the Version 2.2 privacy-domain margins. The initial alternative construction and its post-seed denominator continuation are treated separately, so the privacy statement remains existence-based and local. The results apply only before the earliest admissibility exit or finite privacy-seed boundary. They do not assert global continuation, all-time invariance, post-seed validity, asymptotic residual decay, deadline recovery, active-power sharing, or a simultaneous composite guarantee.
\end{abstract}

\begin{IEEEkeywords}
Islanded AC microgrid, distributed secondary control, public/private coordination, prescribed-performance control, public-history indistinguishability, passive eavesdropper, local-before-exit analysis.
\end{IEEEkeywords}

%======================================================
\section{Introduction}

Islanded AC microgrids rely on distributed secondary control to coordinate inverter-based units while preserving acceptable voltage and frequency behavior. The communication layer is part of this control loop: public coordination states are exchanged among neighboring units, whereas local controller variables and physical measurements may remain private. A useful analysis therefore has to describe both the physical closed loop and the information that is available to an observer of the communication network.

The two requirements are not interchangeable. A controller can be physically regular on an admissible operating region without preventing multiple private explanations of the same communication record. Conversely, a state decomposition that hides local variables does not by itself establish that the coupled microgrid equations remain well posed, that the prescribed-performance coordinates are regular, or that the implemented inputs remain feasible on the design region. The resulting proof problem is a coupled one, but the physical and observation-equivalence arguments must remain logically separate.

This paper analyzes a distributed secondary-control architecture with a public/private virtual coordination layer, a fixed connected cyber graph, and a passive eavesdropper that observes the complete public history. The controller, plant equations, prescribed-performance construction, and privacy equations are treated as frozen design objects. The proof is intentionally local-before-exit: estimates are stated on a selected compact bootstrap/design region and the privacy construction is stopped at the earliest finite-seed or regular-domain boundary. This boundary avoids interpreting a local margin as a global invariance claim.

The paper makes two bounded contributions. First, it establishes a local physical result for the closed loop: local well-posedness, finite compact-dependent component and composite comparison bounds, a finite privacy-residual bound, and actuator feasibility on the verified bootstrap design region. Second, it establishes a local public-history indistinguishability result: after the initial alternative construction, the admissible private realization can be continued while the strict privacy-domain conditions remain valid, and the nominal and alternative realizations generate the same complete passive public history. The two results are reported separately. Deadline recovery, asymptotic residual decay, active-power sharing, forward invariance, global continuation, and a simultaneous composite theorem are outside the final theorem scope.

The contributions of this paper are as follows.
\begin{enumerate}
    \item A local-before-exit analysis of the frozen distributed secondary-control closed loop. On the selected compact bootstrap/design region, the plant, prescribed-performance coordinates, reconstruction interface, and privacy residual are regular, and the composite comparison inequality is available with the stated local feasibility conditions.
    \item A local public-history indistinguishability construction under the passive observation model. The initial non-nominal realization and the subsequent denominator-valid continuation are treated as separate proof obligations, and the claim stops at the finite-seed or regular-domain boundary.
\end{enumerate}


%======================================================
\section{System Model and Problem Formulation}
\label{sec:system_model}

\subsection{Islanded Microgrid and Droop Model}

Consider an islanded AC microgrid containing $N$ inverter-based distributed generators (DGs), indexed by $\mathcal V=\{1,\ldots,N\}$. For DG $i$, let $V_i$, $\omega_i$, and $\delta_i$ denote the voltage amplitude, angular frequency, and phase angle, respectively. The voltage and frequency channel states are
\begin{equation}
\bm{x}_i^V=
\begin{bmatrix}x_{i,1}^V & x_{i,2}^V\end{bmatrix}^{\!T}
=\begin{bmatrix}V_i & \dot V_i\end{bmatrix}^{\!T},
\qquad x_i^\omega=\omega_i,
\qquad \dot\delta_i=\omega_i.
\label{eq:physical_states}
\end{equation}
The reduced droop dynamics are
\begin{align}
\tau_{P_i}\dot\omega_i
&=-(\omega_i-\omega_{\mathrm{ref}})
-k_{P_i}(P_i-P_i^d)-u_i^\omega
+\tau_{P_i}R_i^\omega,
\label{eq:frequency_droop}\\
\tau_{Q_i}k_{V_i}\ddot V_i
&=-(\tau_{Q_i}+k_{V_i})\dot V_i
-(V_i-V_{\mathrm{ref}})
-k_{Q_i}(Q_i-Q_i^d)-u_i^V
+\tau_{Q_i}k_{V_i}R_i^V.
\label{eq:voltage_droop}
\end{align}
Here, $\tau_{P_i}$ and $\tau_{Q_i}$ are filter time constants, $k_{P_i}$ and $k_{Q_i}$ are droop coefficients, $k_{V_i}$ is the voltage-loop coefficient, and $P_i^d$ and $Q_i^d$ are droop setpoints. The terms $R_i^\omega$ and $R_i^V$ collect bounded physical or model uncertainty in their respective channels. Privacy residuals are not included in these terms.

The electrical network is represented by $\mathcal G_e=(\mathcal V,\mathcal E_e)$ with electrical neighbor set $\mathcal N_i^e$. Under the frozen lossless inductive specialization, let $B_{ik}$ denote the line susceptance magnitude and define $\delta_{ik}=\delta_i-\delta_k$. The active and reactive powers are
\begin{align}
P_i
&=P_i^L(V_i)+\sum_{k\in\mathcal N_i^e}
V_iV_k|B_{ik}|\sin(\delta_{ik}),
\label{eq:active_power_flow}\\
Q_i
&=Q_i^L(V_i)+V_i^2\sum_{k\in\mathcal N_i^e}|B_{ik}|
-\sum_{k\in\mathcal N_i^e}V_iV_k|B_{ik}|\cos(\delta_{ik}),
\label{eq:reactive_power_flow}
\end{align}
where $P_i^L(V_i)$ and $Q_i^L(V_i)$ are the local voltage-dependent load functions. Conductance terms are not absorbed into $R_i^\omega$ or $R_i^V$.

For later analysis, define the known channel drifts
\begin{align}
F_i^\omega
&=-\frac{(\omega_i-\omega_{\mathrm{ref}})
+k_{P_i}(P_i-P_i^d)}{\tau_{P_i}},
\label{eq:frequency_drift}\\
F_i^V
&=-\frac{(\tau_{Q_i}+k_{V_i})x_{i,2}^V
+(x_{i,1}^V-V_{\mathrm{ref}})
+k_{Q_i}(Q_i-Q_i^d)}{\tau_{Q_i}k_{V_i}}.
\label{eq:voltage_drift}
\end{align}
Equations~\eqref{eq:frequency_droop} and \eqref{eq:voltage_droop} can then be written as
\begin{align}
\dot x_{i,1}^V&=x_{i,2}^V,\nonumber\\
\dot x_{i,2}^V&=F_i^V-\frac{u_i^V}{\tau_{Q_i}k_{V_i}}+R_i^V,
\label{eq:voltage_input_affine}\\
\dot\omega_i&=F_i^\omega-\frac{u_i^\omega}{\tau_{P_i}}+R_i^\omega.
\label{eq:frequency_input_affine}
\end{align}
The privacy-aware coordination layer reaches the plant through the unique interface
\begin{equation}
u_i^V=\widehat c_i^V,
\qquad
u_i^\omega=\widehat c_i^\omega.
\label{eq:plant_interface}
\end{equation}
No public or private coordination state enters \eqref{eq:voltage_input_affine}--\eqref{eq:frequency_input_affine} through another path.

\subsection{Electrical and Cyber Graphs}

The electrical graph $\mathcal G_e$ determines the power-flow coupling in \eqref{eq:active_power_flow}--\eqref{eq:reactive_power_flow}. Communication is described separately by a fixed, connected, undirected cyber graph $\mathcal G_c=(\mathcal V,\mathcal E_c)$. Its symmetric adjacency weights and Laplacian are
\begin{equation}
a_{ij}=a_{ji}>0\quad\text{if }(i,j)\in\mathcal E_c,
\qquad
L_c=\operatorname{diag}\!\left(\sum_j a_{ij}\right)-[a_{ij}],
\label{eq:cyber_graph}
\end{equation}
with $a_{ij}=0$ otherwise. Reference access is encoded by $b_i\geq0$, with at least one pinned DG in each regulated channel. The corresponding pinned matrix is $L_c+\operatorname{diag}(b_i)$. The electrical and cyber neighbor sets, $\mathcal N_i^e$ and $\mathcal N_i^c$, are distinct model objects and need not coincide.

\subsection{Public/Private Coordination Interface}

For each channel $\nu\in\{V,\omega\}$, $p_i^\nu$ denotes the public coordination substate and $q_i^\nu$ denotes the private substate retained by DG $i$. The private weights, residual variables, and locally reconstructed commands are also retained locally. The regular message transmitted by DG $i$ is exactly
\begin{equation}
\bm{m}_i(t)=
\begin{bmatrix}p_i^V(t)&p_i^\omega(t)\end{bmatrix}^{T}.
\label{eq:public_message}
\end{equation}
Thus, raw $V_i$, $\omega_i$, $P_i$, $Q_i$, the nominal coordination state $c_i$, the reconstructed command $\widehat c_i$, $q_i$, privacy residuals, private weights, physical uncertainty, and private controller memory are not regular message components. This ownership convention does not mask the physical state from an adversary that obtains an independent physical-sensor channel; such a channel is outside the frozen observation model.

\subsection{Passive Observation Map}

At time $t$, DG $i$ has access to the neighboring message history
\begin{equation}
\{\bm{m}_j(s):j\in\mathcal N_i^c,\ 0\leq s\leq t\}.
\label{eq:neighbor_history}
\end{equation}
The passive eavesdropper observes the complete public history
\begin{align}
\mathcal O_{\mathrm{adv}}[0,t]
={}&\{\bm{m}_j(s):j\in\mathcal V,\ 0\leq s\leq t\}
\cup \mathcal H_c[0,t]\nonumber\\
&\cup\{\mathcal G_c,a_{ij},b_i,V_{\mathrm{ref}},
\omega_{\mathrm{ref}},\rho_i^V(\cdot),\rho_i^\omega(\cdot),
T_V,T_\omega,\text{ declared public controller parameters}\}.
\label{eq:adversary_history}
\end{align}
Here, $\mathcal H_c[0,t]$ contains the disclosed topology, timing, and protocol metadata. Public schedules and their parameters are included because they are part of the disclosed protocol, not because a deadline property is asserted. Private substates, private weights, residual and internal privacy variables, private controller memory, raw physical-sensor histories, and locally reconstructed commands are not observed unless separately declared public. The adversary records but does not alter public messages; active manipulation and communication failures are outside this model.

Two admissible private initializations are publicly indistinguishable on a declared interval when they generate the same $\mathcal O_{\mathrm{adv}}[0,t]$ under the same public references, cyber graph, schedules, and protocol metadata while their protected initial local virtual coordination states differ. This definition concerns non-unique reconstruction from the complete public history. It does not describe differential privacy or cryptographic secrecy.

\subsection{Local-Before-Exit Problem Statement}

The physical problem is to characterize the frozen closed loop from an admissible independent initial state while the solution remains in the open regular domain $\mathcal D_{\min}$ and, for compact-dependent estimates, in a selected compact bootstrap/design region $\mathcal K_0\Subset\mathcal D_{\min}$. The requested result is local well-posedness together with the finite component and composite comparison bounds available before the first admissibility exit. The problem does not require exclusion of that exit.

The privacy problem is to determine whether a non-nominal admissible private realization can produce the same complete passive public history as a nominal realization. The construction is restricted to the Version~2.2 schedule-regular privacy domain and stops at the earliest finite-seed or regular-domain boundary. No conclusion is sought after that stopping boundary. Section~III formalizes the active assumptions and domains used by these two local problems.

\section{Definitions and Active Assumptions}
\label{sec:definitions_assumptions}

\subsection{Independent and Reconstructed Coordinates}

The local solution is formulated on the independent coordinates
\begin{equation}
X=\operatorname{col}_{i\in\mathcal V}\!\left(V_i,\dot V_i,\omega_i,\delta_i,
p_i^V,q_i^V,p_i^\omega,q_i^\omega\right).
\label{eq:independent_coordinates}
\end{equation}
The physical errors, normalized errors, transformed errors, virtual controls,
reconstructed commands, and privacy residuals are algebraic or smooth images of
$(X,t)$ through the frozen maps. In particular, $e_{i,0}^V$, $e_{i,0}^\omega$,
$\sigma_i^V$, $\sigma_i^\omega$, $\zeta_i^V$, $\zeta_i^\omega$, $r_i^V$,
and $r_i^\omega$ are reconstructed analysis quantities; they are not additional
independent states. The augmented bookkeeping vector used in the proof is
therefore interpreted as the image of \eqref{eq:independent_coordinates}, with
its consistency derivatives supplied by the frozen equations.

\subsection{Admissible Open Domain and Bootstrap Region}

Let $\mathcal D_{\min}$ denote the genuine open regular domain of the independent
coordinates on which the plant, controller, prescribed-performance maps, and
privacy divisions are defined. It includes strict funnel inequalities
$|\sigma_i^V|<1$ and $|\sigma_i^\omega|<1$, positive performance envelopes,
the regular privacy denominators, and the physical operating domain inherited
from the frozen model. A compact set
\begin{equation}
\mathcal K_0\Subset\mathcal D_{\min}
\label{eq:bootstrap_region}
\end{equation}
is selected with the admissible initial condition in its interior. The set
$\mathcal K_0$ is a design and estimation region for local bounds; it is not
declared forward invariant. The first-exit time is the earliest time at which
the independent trajectory leaves $\mathcal D_{\min}$ or a separately declared
local design margin ceases to hold.

\subsection{Definition 1: Closed-Loop Local Solution}

\textbf{Definition 1 (closed-loop local solution).}
Given an admissible initial condition in $\mathcal D_{\min}$, a closed-loop local
solution is an absolutely continuous independent trajectory $X(t)$ satisfying
the frozen plant, graph, controller, and public/private coordination equations
almost everywhere on a maximal interval $[0,\tau_{\mathrm{ex}})$, where
$\tau_{\mathrm{ex}}$ is the first admissibility-exit time. All reconstructed
coordinates are evaluated from $X(t)$ and remain consistent with their defining
maps on this interval. This definition asserts local well-posedness and a
stopping boundary only; it does not assert continuation beyond
$\tau_{\mathrm{ex}}$.

\subsection{Definition 2: Public-History Indistinguishability}

\textbf{Definition 2 (public-history indistinguishability).}
Two admissible realizations are publicly history-indistinguishable on an interval
$[0,\tau)$ if their complete passive observations satisfy
\begin{equation}
\mathcal O_{\mathrm{adv}}[0,t]=\mathcal O'_{\mathrm{adv}}[0,t],
\qquad 0\leq t<\tau,
\label{eq:definition_public_history}
\end{equation}
including every transmitted public message and every disclosed element of
$\mathcal H_c[0,t]$, while at least one protected private initial quantity
$S_i$ differs from its alternative $S_i'$. The observation is passive and
complete with respect to the declared public payload and metadata. Private
memory, private weights, residuals, locally reconstructed commands, and raw
physical-sensor histories are excluded unless explicitly declared public. The
definition is existence-based and does not assert cryptographic secrecy or
uniqueness failure outside the declared local stopping interval.

\subsection{Assumption 1: Local Physical and Graph Regularity}

\textbf{Assumption 1.} The frozen physical model is defined on the declared
operating region $\Delta$, its channel parameters and positive time constants
are regular, and the load and power-flow maps are locally regular there. The
cyber graph $\mathcal G_c$ is fixed, connected, undirected, and has the declared
pinning needed by the distributed errors; the electrical graph remains a
separate plant object. References and performance envelopes are public and
regular with strictly positive radii. The uncertainty signals $R_i^V$ and
$R_i^\omega$ are measurable and locally essentially bounded on compact subsets
of $\Delta$ with the declared bounds. The initial condition satisfies the
strict funnel feasibility conditions and belongs to the interior of
$\mathcal D_{\min}$. The feasible actuator sets $U_i$ and frozen input maps are
locally regular; actuator feasibility on a selected $\mathcal K_0$ is a design
check performed by PO-13, not a global premise of this assumption.

\subsection{Assumption 2: Version 2.2 Privacy-Domain Regularity}

\textbf{Assumption 2.} For every affected agent/channel pair, the nominal
initial private split has the declared nonzero margin $\eta_{z,i}^\nu>0$,
the nominal private weights lie strictly inside their ES-46 bounds with margin
$\eta_{w,i}^\nu>0$, and the prescribed privacy schedule satisfies
$\gamma_{\mathrm{priv},i}^\nu(t)\geq\eta_{\gamma,i}^\nu>0$ on the common finite
seed interval $I_s=[0,T_s]$. The margins satisfy the Version 2.2 feasibility
relations, and the privacy gains and finite local residual budget obey the
declared Privacy Gain Feasibility Condition on the selected compact design
region. The schedule is regular enough for the finite local residual estimate;
no asymptotic residual conclusion is included here. The adversary is passive
and observes exactly the complete public history in Definition 2.

Assumption 2 does not assert the existence of a non-nominal admissible
alternative, a positive perturbation radius, or a compatible private trajectory:
those are conclusions of PO-04. It also does not assert denominator validity
beyond the initial construction interval: any additional ES-60--ES-61 validity,
continuation, or compatible extension up to the retained privacy stopping
boundary remains the conclusion of PO-05. Neither assumption asserts global
invariance or continuation.

\section{Local Physical Analysis}
\label{sec:local_physical_analysis}

\subsection{Prescribed-Performance Coordinates and Frozen Controller}

For $\nu\in\{V,\omega\}$, the frozen performance envelope is
\begin{equation}
\rho_i^\nu(t)
=\rho_{i,\infty}^\nu+
  (\rho_{i,0}^\nu-\rho_{i,\infty}^\nu)h(t/T_\nu),
\label{eq:ppc_envelope}
\end{equation}
\begin{equation}
h(s)=
\begin{cases}
1-10s^3+15s^4-6s^5,&0\leq s\leq1,\\
0,&s>1.
\end{cases}
\label{eq:ppc_shape}
\end{equation}
The normalized and transformed coordinates are
\begin{equation}
\sigma_i^\nu=\frac{e_{i,0}^\nu}{\rho_i^\nu},
\qquad
\zeta_i^\nu=\operatorname{atanh}(\sigma_i^\nu),
\qquad
h_i^\nu=\frac{1}{\rho_i^\nu[1-(\sigma_i^\nu)^2]}.
\label{eq:ppc_coordinates}
\end{equation}
They are well defined only in the strict interior $|\sigma_i^\nu|<1$.
For the voltage channel, the retained virtual control and backstepping error are
\begin{align}
\alpha_i^V
&=\sigma_i^V\dot\rho_i^V
-k_1^V\rho_i^V[1-(\sigma_i^V)^2]\zeta_i^V,
\nonumber\\
\chi_i^V&=x_{i,2}^V-\alpha_i^V,
\label{eq:voltage_virtual_control}
\end{align}
whereas the frequency channel uses
\begin{equation}
\alpha_i^\omega
=\sigma_i^\omega\dot\rho_i^\omega
-k_1^\omega\rho_i^\omega[1-(\sigma_i^\omega)^2]\zeta_i^\omega.
\label{eq:frequency_virtual_control}
\end{equation}
The frozen nominal commands and the privacy-aware plant inputs are
\begin{align}
c_i^V
&=\tau_{Q_i}k_{V_i}\!\left(F_i^V-\dot\alpha_i^V+k_2^V\chi_i^V
+h_i^V\zeta_i^V\right)+k_c^V e_i^V,
\label{eq:local_voltage_command}\\
c_i^\omega
&=\tau_{P_i}(F_i^\omega-\alpha_i^\omega)+k_c^\omega e_i^\omega,
\label{eq:local_frequency_command}\\
u_i^V&=c_i^V+r_i^V,
\qquad
u_i^\omega=c_i^\omega+r_i^\omega.
\label{eq:local_reconstructed_inputs}
\end{align}
Substitution through the unique plant interface gives the frozen transformed
subsystems
\begin{align}
\dot\zeta_i^V
&=-k_1^V\zeta_i^V+h_i^V\chi_i^V,
\nonumber\\
\dot\chi_i^V
&=-k_2^V\chi_i^V-h_i^V\zeta_i^V
-\frac{k_c^V e_i^V+r_i^V}{\tau_{Q_i}k_{V_i}}+R_i^V,
\label{eq:local_voltage_subsystem}\\
\dot\zeta_i^\omega
&=-k_1^\omega\zeta_i^\omega
+h_i^\omega\!\left(
-\frac{k_c^\omega e_i^\omega+r_i^\omega}{\tau_{P_i}}
+R_i^\omega\right).
\label{eq:local_frequency_subsystem}
\end{align}
These identities retain the privacy residual explicitly and do not absorb it
into the physical uncertainty.

\subsection{Local Well-Posedness and First-Exit Interval}

On $\mathcal D_{\min}$, the load and power-flow maps, graph maps, PPC
transformations, and algebraic controller maps are locally regular. The product
$g_i^\nu z_i^\nu$ in the privacy dynamics is locally Lipschitz in the private
difference, including at $|z_i^\nu|=\gamma_{\mathrm{priv},i}^\nu$.
Together with the measurable, locally essentially bounded uncertainty and
privacy coefficients in Assumptions 1--2, the reduced independent vector field
satisfies the Caratheodory conditions. Hence PO-16A supplies a unique absolutely
continuous maximal local solution from every admissible initial condition.

All analysis below is restricted to compact time intervals before
$\tau_{\mathrm{ex}}$, while the independent trajectory remains in
$\mathcal D_{\min}$ and in the selected $\mathcal K_0$. On such intervals, all
reconstructed coordinates and compact-dependent constants are finite. This
local result neither excludes a later admissibility exit nor extends the
trajectory beyond it.

\subsection{Command-Rate and Finite Residual Bounds}

For each channel, subtraction and averaging of the frozen public/private
dynamics give
\begin{align}
\dot z_i^\nu
&=-[\lambda_{\mathrm{tr},i}^\nu+w_{i,12}^\nu
+w_{i,21}^\nu g_i^\nu]z_i^\nu,
\label{eq:local_z_dynamics}\\
\dot r_i^\nu
&=-\lambda_{\mathrm{tr},i}^\nu r_i^\nu
+\frac{1}{2}(w_{i,12}^\nu-w_{i,21}^\nu g_i^\nu)z_i^\nu
-\dot c_i^\nu.
\label{eq:local_r_dynamics}
\end{align}
PO-01 gives the channel-wise exponential bound on $z_i^\nu$. PO-03 then gives
finite compact-dependent command-rate constants
\begin{equation}
\|\dot{\bm c}^V\|\leq C_c^V(\mathcal K_0),
\qquad
\|\dot{\bm c}^\omega\|\leq C_c^\omega(\mathcal K_0),
\label{eq:local_command_rate_bounds}
\end{equation}
before the first exit while the trajectory remains in $\mathcal K_0$.
Since $r_i^\nu(0)=0$, variation of constants in
\eqref{eq:local_r_dynamics} yields
\begin{align}
r_i^\nu(t)=\int_0^t &e^{-\lambda_{\mathrm{tr},i}^\nu(t-s)}
\bigg[\frac{1}{2}(w_{i,12}^\nu-w_{i,21}^\nu g_i^\nu)z_i^\nu
\nonumber\\[-1mm]
&\hspace{29mm}-\dot c_i^\nu\bigg](s)\,ds.
\label{eq:local_residual_convolution}
\end{align}
PO-02A therefore supplies a finite residual bound on the same compact local
interval. A bounded command rate supports only this finite convolution estimate;
no residual convergence conclusion follows from it.

\subsection{Voltage and Frequency Component Inequalities}

Use only the frozen Lyapunov components
\begin{align}
\mathscr V_V
&=\frac{1}{2}\sum_i\!\left[p_\zeta^V(\zeta_i^V)^2
+p_\chi^V(\chi_i^V)^2\right],
\nonumber\\
\mathscr V_\omega
&=\frac{1}{2}\sum_i(\zeta_i^\omega)^2,
\nonumber\\
\mathscr V_{\mathrm{priv}}
&=\frac{1}{2}\sum_{i,\nu}p_c^\nu
\left[(z_i^\nu)^2+(r_i^\nu)^2\right],
\nonumber\\
\mathscr V_{\mathrm{cl}}
&=\mathscr V_V+\mathscr V_\omega+\mathscr V_{\mathrm{priv}}.
\label{eq:frozen_lyapunov_components}
\end{align}
On $\mathcal K_0$, define the frozen disturbance aggregates
\begin{align}
D_i^V
&=\frac{k_c^V|e_i^V|+|r_i^V|}{\tau_{Q_i}k_{V_i}}
+\bar R_i^V,
\nonumber\\
D_i^\omega
&=\frac{k_c^\omega|e_i^\omega|+|r_i^\omega|}{\tau_{P_i}}
+\bar R_i^\omega.
\label{eq:component_disturbances}
\end{align}
PO-08 and PO-09 establish, pointwise on the selected compact region and before
the first admissibility exit,
\begin{align}
\dot{\mathscr V}_V
&\leq-a_{Vz}\|\bm\zeta^V\|^2-a_{V\chi}\|\bm\chi^V\|^2
+\sum_i\frac{(p_\chi^V)^2(D_i^V)^2}{2\epsilon_{V2}},
\label{eq:local_voltage_inequality}\\
\dot{\mathscr V}_\omega
&\leq-a_{\omega z}\|\bm\zeta^\omega\|^2
+\sum_i\frac{(\bar h_i^\omega)^2(D_i^\omega)^2}
{2\epsilon_\omega},
\label{eq:local_frequency_inequality}
\end{align}
where $a_{Vz}$, $a_{V\chi}$, and $a_{\omega z}$ are the positive coefficients
fixed by ES-94, ES-95, and ES-98. Under the Privacy Gain Feasibility Condition,
PO-10 gives on the same local region
\begin{equation}
\dot{\mathscr V}_{\mathrm{priv}}
\leq-a_z\|\bm z\|^2-a_r\|\bm r\|^2
+d_c\|\dot{\bm c}\|^2,
\label{eq:local_privacy_inequality}
\end{equation}
with $a_z>0$, $a_r>0$, and finite $d_c$. This component estimate uses
\eqref{eq:local_command_rate_bounds}, not a residual-decay premise.

\subsection{Bootstrap Actuator/Funnel Feasibility}

PO-13 checks feasibility on $\mathcal K_0$ before the composite gain closure.
For the frozen reconstructed commands, let
\begin{equation}
\bar U_i^\nu(\mathcal K_0)
=\sup_{X\in\mathcal K_0}|\widehat c_i^\nu(X)|.
\label{eq:local_actuator_demand}
\end{equation}
The voltage and frequency command maps give the finite bounds
\begin{align}
\bar U_i^V
&\leq\tau_{Q_i}k_{V_i}
(\bar F_i^V+\overline{\dot\alpha}_i^V
+k_2^V\bar\chi_i^V+\bar h_i^V\bar\zeta_i^V)
\nonumber\\
&\quad+|k_c^V|\bar e_i^V+\bar r_i^V,
\label{eq:voltage_actuator_bound}\\
\bar U_i^\omega
&\leq\tau_{P_i}(\bar F_i^\omega+\bar\alpha_i^\omega)
+|k_c^\omega|\bar e_i^\omega+\bar r_i^\omega.
\label{eq:frequency_actuator_bound}
\end{align}
The symbolic design test requires the resulting compact command intervals to be
strictly contained in the declared actuator sets $U_i^\nu$, together with the
frozen funnel, gain, and privacy-feasibility inequalities. It certifies actuator
and funnel feasibility only on the selected design region while the trajectory
remains there. It does not establish persistence of those margins after an exit.

\subsection{Composite Local Comparison}

The graph identity used by PO-06 is
\begin{equation}
\bm e=\operatorname{diag}(b_i)\bm e_0+L_c\bm c
+L_c\bm r+\frac{1}{2}L_c\bm z.
\label{eq:local_graph_identity}
\end{equation}
Together with the frozen algebraic command maps, it bounds $\bm e$ and $\bm c$
by the physical, residual, and decomposition coordinates without treating the
graph term as an external input. The component inequalities are assembled with
the closed PO-07 certificate
\begin{equation}
Q_{\mathrm{cl}}=Q_0-H^T W_DH\succ0,
\qquad
\lambda_{\mathrm{cl}}=\lambda_{\min}(Q_{\mathrm{cl}})>0,
\label{eq:local_composite_certificate}
\end{equation}
where the frequency factor $(\bar h_i^\omega)^2$ occurs exactly once in
$W_D^\omega$. Consequently, on the selected $\mathcal K_0$, before the first
admissibility exit and while the solution remains in $\mathcal D_{\min}$,
\begin{equation}
\dot{\mathscr V}_{\mathrm{cl}}
\leq-a_{\mathrm{cl}}\mathscr V_{\mathrm{cl}}+d_R+d_{\mathrm{priv}}(t),
\qquad a_{\mathrm{cl}}>0,
\label{eq:local_composite_inequality}
\end{equation}
and comparison gives
\begin{align}
\mathscr V_{\mathrm{cl}}(t)
\leq{}&e^{-a_{\mathrm{cl}}t}\mathscr V_{\mathrm{cl}}(0)
\nonumber\\
&+\int_0^t e^{-a_{\mathrm{cl}}(t-s)}
[d_R+d_{\mathrm{priv}}(s)]\,ds.
\label{eq:local_comparison_bound}
\end{align}
Here $d_R$ contains the declared physical-uncertainty and compact-region terms,
whereas $d_{\mathrm{priv}}$ is finite by PO-03 and PO-02A. Equations
\eqref{eq:local_composite_inequality}--\eqref{eq:local_comparison_bound} are
local comparison statements. They do not exclude a funnel, actuator, physical,
or loss-of-compactness exit and do not provide a conclusion after such an exit.


%======================================================
\section{Illustrative Simulation}

\section{Hardware-in-the-Loop Closed-Loop Test}

\section{Conclusion}



%======================================================
\begin{thebibliography}{99}
\bibitem{ref1}
Author. Title. Journal, Year.

\end{thebibliography}
\end{document}
